% =============================================================================
% pre-textuais/resumo.tex  --  Resumo em português e abstract em inglês
% =============================================================================

% ----- Resumo em português -----
\begin{resumoumacoluna}
A colônia de gatos comunitários do Campus~2 da Universidade de São Paulo em São Carlos é acompanhada pela Atividade de Extensão Gatos do C2, que opera dez pontos de alimentação distribuídos pelo campus. Como apoio ao manejo dessa colônia, especifica-se uma arquitetura modular de aquisição de imagem baseada em câmeras IP comerciais em regime \textit{offline-first}, sobre a qual se organiza um \textit{pipeline} de visão computacional em quatro fases: triagem por movimento, detecção de fauna, classificação de espécie e reidentificação individual (re-ID). O \textit{pipeline} apoia-se em modelos abertos pré-treinados (MegaDetector, SpeciesNet, MegaDescriptor e PetFace) e foi avaliado em ambiente \textit{single-machine} com GPU NVIDIA MX250. A cascata de detecção e classificação mostra-se aplicável a fontes heterogêneas, com convergência dependente da origem dos dados. No estágio de re-ID, o Top-1 em \textit{closed-set} reduz-se de 0,720 para 0,590 com o aumento do catálogo de $N{=}50$ para $N{=}200$ identidades; em \textit{open-set}, a AUC-ROC permanece próxima de 0,65 e a TPR@FPR=1\% recua de 0,213 para 0,087. Os resultados delimitam o uso de modelos abertos sem ajuste fino em colônias urbanas e indicam, como continuidade, a coleta supervisionada de dados em campo, o ajuste fino sobre o catálogo individual da colônia e a integração com plataforma de gestão.

\vspace{\onelineskip}
\noindent
\textbf{Palavras-chave:} visão computacional, câmeras IP, reidentificação animal, monitoramento de fauna, gatos comunitários, USP São Carlos, extensão universitária.
\end{resumoumacoluna}

% ----- Abstract em inglês -----
\begin{otherlanguage*}{english}
\renewcommand{\resumoname}{ABSTRACT}
\begin{resumoumacoluna}
The community-cat colony at Campus~2 of the University of São Paulo at São Carlos is supported by the Gatos do C2 university extension activity, which operates ten feeding points across the campus. As support to the management of this colony, a modular image-acquisition architecture based on commercial IP cameras under an \textit{offline-first} regime is specified, on which a four-stage computer-vision \textit{pipeline} is organised: motion-based triage, fauna detection, species classification and individual re-identification (re-ID). The \textit{pipeline} relies on open pre-trained models (MegaDetector, SpeciesNet, MegaDescriptor and PetFace) and was evaluated on a \textit{single-machine} environment with an NVIDIA MX250 GPU. The detection-and-classification cascade is applicable to heterogeneous sources, with source-dependent convergence. At the re-ID stage, closed-set Top-1 declines from 0.720 to 0.590 as the catalogue grows from $N{=}50$ to $N{=}200$ identities; under open-set conditions, AUC-ROC remains close to 0.65 and TPR@FPR=1\% recedes from 0.213 to 0.087. The results delimit the use of open models without fine-tuning in urban colonies and point, as continuity, to supervised field-data collection, fine-tuning on the colony individual catalogue and integration with a colony management platform.

\vspace{\onelineskip}
\noindent
\textbf{Keywords:} computer vision, IP cameras, animal re-identification, wildlife monitoring, community cats, USP São Carlos, university extension.
\end{resumoumacoluna}
\end{otherlanguage*}